% !TeX root = ../papers/nc-p3.tex
\subsection{平面グラフ基底の固定直線と指定点を忘れた軌道}
\label{sec:plane-graph-ore-windows}
命題\ref{prop:ore-mark-recognition}の三生成基底の点対応が射影平面全体の自己同型
$[G]$ のグラフである場合を考える。三生成関係は $p\wedge G^{-1}r$ の三成分で、
$B$ は多項式環の次数付きtwistである。実際、長さ$n$のテンソルを
$(p,Gp,\ldots,G^{n-1}p)$ で評価すると、各位置を順次 $G$ の冪で
基底変更した後の関係は可換関係となり、商は全$n$次式である。

\begin{theorem}\label{thm:plane-ore-fixed-lines}
平面グラフについて、指定列を持つOre補正の空間 $P_B$ と移動 $F_B$ は次で尽くされる。
\begin{enumerate}
\item $[G]=1$ なら二次導分をEuler導分の一次式倍で割った十五次元空間で、移動は恒等である。
\item $[G]\ne1$ が直線を点ごとに固定しなければ $P_B=0$ である。
\item 残りでは $G=1+e\ell$ と取る。ここで $e$ は点ベクトル、$\ell$ は一次形式で、
$t=1+\ell(e)\ne0$ とする。$P_B$ は固定直線 $\ell=0$ 上の二次式の三次元空間、
移動は $t$ 倍である。全ての例の最小周期は1である。
\end{enumerate}
\end{theorem}
\begin{proof}
恒等の場合、三次の条件は $(d_{i+1}-d_i)(h)k=(d_{i+1}-d_i)(k)h$ である。
従って差が一次式倍のEuler導分に限り、定理\ref{thm:ternary-ore-linear-space}から結論を得る。
非恒等の場合、定理\ref{thm:nodal-tangent-windows}の平面における二つの縮約平面の計算により
$d_i$ の評価は一般点で $p,Gp$ の一次結合となる。この二項を
$a_i(p)Gp-b_{i+1}(Gp)p$ と書く。斉次性から係数は $\mathcal O(1)$ の有理切断である。
混合三次関係を $(p,Gp,G^2p)$ で評価すると
$(a_{i+1}-b_{i+1})(Gp)$ と三関係の $(p,G^2p)$ での評価の積が零である。
後者が全て消えるのは $Gp=p$ が射影的に成立するときなので、$a_i=b_i$ である。

二次小行列で二係数を解くと、極因子は固定直線に限る。実際、固定点集合は
固有空間の射影化の和で、平面上の有理切断は余次元一の全点で正則なら正則である。
従って固定直線がなければ補生成元を変えて全補正を消せる。
固定直線があればその上での固有値を1に選ぶと $G=1+e\ell$ となり、
$p\wedge Gp=\ell(p)p\wedge e$ のため一般の極は高々単純である。
よって一意な主要部は $a_i=f_i/\ell$ の二次分子の直線への制限で与えられる。
$\ell(Gp)=t\ell(p)$ より二次補正の正則性は
$f_{i+1}|_{\ell=0}=t f_i|_{\ell=0}$ と同値である。逆に任意の二次式 $f$ に対し
$a_i=t^i f/\ell$ の表示は全位置で正則な二次補正を与える。
有理式での相殺と三次評価の単射性により全混合三次関係を満たす。
正則な一次切断の加算はちょうど補生成元変更である。初項が左右乗法の和なら
一般点での係数の一意性から主要部は零であり、主張の同定を得る。
\end{proof}

\begin{theorem}\label{thm:plane-ore-unmarked}
非恒等の平面グラフ基底を少なくとも一つ持つOre窓どうしの添字保持同型を考える。
固定直線のない型は $[G]$ の射影共役類で分類される。
固定直線のある型は次の表で分類され、表の各行および固定直線のない型は互いに素である。
$r$ は二元二次式 $f|_{\ell=0}$ の階数である。
\[
\begin{array}{c|l}
 t\ne1 & \{t,t^{-1}\},\quad r=0,1,2\\
 t=1 & r=0;\quad r=1,2\ \text{ではさらに}\ f(e)=0\ \text{か否か}
\end{array}
\]
全係数で正則である。各 $t=1$ の類は $t\ne1$ の類の平坦特殊化である。
\end{theorem}
\begin{proof}
固定直線のある型に対し、位置 $i$ の旧水平座標を $G^ip$ で置換する。
旧第四座標を $u+t^i f(G^ip)/\ell(G^ip)-f(p)/\ell(p)$ とする。
後二項の差は一次式である。実際その分子は直線上で消え、$G^i$ は直線上恒等である。
新しい座標 $(p,u),(r,v)$ では六関係は
\[
 p\wedge r=0,\qquad u G r-vp=e f(p,r)
\]
となる。$f(p,r)$ は $f$ の対称偏極である。これは純粋な三関係を移し、
$r=p$ で旧有理補正を代入して $G^{-i}$ を掛ければ確認できる。
対角で消える双一次式は純粋な可換関係に吸収される。
特に一定表示である。$C=k[p_1,p_2,p_3]$ に対する通常のOre拡大としても、
$\gamma(h)(p)=h(G^{-1}p)$、$\delta(h)=h(G^{-1}e)f$ である。
$h-\gamma(h)$ が $h(e)\ell/t$ に等しいため、$\delta$ の捩れたLeibniz則は可換関係を保つ。
従って順序単項式基底があり、また定理\ref{thm:maximal-tangent-ore}で全次数の正則性を得る。

$Q$ を二次曲面 $u\ell(p)-f(p)=0$、$D$ を直線 $p\wedge e=0$ とする。
辺の第一射影はイデアル $\mathcal I_Q\mathcal I_D$ のスキームへ同型である。
$p\ne0$ の開集合では $r=p$ と取れ、混合式の一行で $v$ を消すと
$(u\ell(p)-f(p))p\wedge e=0$ が残る。$p=0$ の近傍では $u=v=1$ と取れ、
$(G-e f(p,-))r=p$ である。左の係数はそこで可逆で、階数一更新の余因子を掛けると
$p\wedge r$ は単元倍を除いて同じ積となる。辺の対応は $Q$ 上恒等である。
$f$ を $\ell$ と独立な二変数で取れば $\operatorname{rk}Q=2+r$ であり、
その被約な曲面部分から $Q$、イデアル商から $D$ を復元できる。
各位置の射影同型は $Q$ 上で隣の位置と等しい。$Q$ を点ごとに固定する射影変換は恒等である。
実際、固定集合は固有空間の和で、平面を二枚含む場合も二固有値は交線上で一致する。
従って一つの射影同型で全位置を比較できる。

$t\ne1$ のとき $e$ を第三方向に選び、$\ell=(t-1)z$、$f$ を前二変数の式に取れる。
$Q\cap D$ は異なる二点で、$D$ 上の対応は $u/z\mapsto t u/z$ である。
従って倍率は反転を除いて不変である。前二変数の基底変更と倍率で同階数の $f$ を移せる。
$u$ と $z$ を交換すると最後の関係 $tuz'-zv=f(p,r)$ は倍率変更で逆倍率の表示となり、
非零な $f$ の定数倍も吸収できる。$t=1$ なら $e\in\ker\ell$ で、
$Q$ は $D$ を含むか、$D$ と一点で二重に交わる。この二つを区別する条件は $f(e)=0$ である。
逆に固定直線上で $[e]$ を保つ二変数の射影変換は全て $[G]$ と可換な変換に延びる。
実際 $e$ の倍率に等しく $\ell$ の商方向を伸縮すればよい。
この群は射影直線の $[e]$ の補集合上のアフィン群であるから、
二元二次式の根の重複度と $[e]$ が根に含まれるかで軌道が定まる。

固定直線がないときは $f=0$、同じ位置変更で $p\wedge r=0,\ u G r=vp$ となる。
辺の台の二次元成分は平面 $u=0$ 一枚だけで、それ以外は固有点と指定点を結ぶ有限本の直線である。
これで二次曲面を持つ型から区別できる。この平面の外では $u=v=1$ で $r=G^{-1}p$、
接空間が三次元の点は $p=0$ だけである。実際 $G^{-1}p=\mu p$ となる非零点では
微分は $p\wedge(G^{-1}-\mu)dp$ で、その恒等消滅は $G^{-1}-\mu$ の階数が高々1となることを強いる。
従って同型は指定点を保ち、指定同型の基底共役条件が必要十分である。

最後に $e=(1,0,0)$、$\ell=z+s x$、$f=f(x,y)$ と取ると $t=1+s$ である。
$s=0$ で所定の任意の移行型を得る。$1+s$ を逆にした係数環上でも上記Ore拡大は
順序単項式基底を持つので、この特殊化は平坦である。固定した $G$ で $f$ の任意の退化も同様である。
記号的な余因子・三次重なりの検算と、Jordan各型の有理標本での独立した
補正空間検算は \path{computations/verify_plane_graph_ore.py} にある。
\end{proof}
